%% LyX 2.4.4 created this file.  For more info, see https://www.lyx.org/.
%% Do not edit unless you really know what you are doing.
\documentclass[english]{article}
\usepackage[T1]{fontenc}
\usepackage[latin9]{inputenc}
\usepackage{enumitem}
\usepackage{amsmath}
\usepackage{amsthm}
\usepackage{amssymb}
\usepackage{geometry}
\geometry{verbose,tmargin=1in,bmargin=1in,lmargin=1in,rmargin=1in}

\makeatletter
%%%%%%%%%%%%%%%%%%%%%%%%%%%%%% Textclass specific LaTeX commands.
\numberwithin{equation}{section}
\numberwithin{figure}{section}
\newlength{\lyxlabelwidth}      % auxiliary length 
\theoremstyle{plain}
\newtheorem{thm}{\protect\theoremname}[section]
\theoremstyle{definition}
\newtheorem{xca}[thm]{\protect\exercisename}

%%%%%%%%%%%%%%%%%%%%%%%%%%%%%% User specified LaTeX commands.
\usepackage{fancyhdr}
\usepackage{lastpage}
\pagestyle{fancy}
\usepackage{enumitem}
\usepackage{ifthen}
\everymath=\expandafter{\the\everymath\displaystyle}
%\DeclareMathSizes{10}{10}{10}{10}
\usepackage{multicol}

\usepackage{tikz}
\usetikzlibrary{patterns}
\usetikzlibrary{plotmarks}
\usepackage{pgfplots}
\pgfplotsset{compat=newest}

\usepackage{calc}
\usepackage[nomessages]{fp}

\usepackage{datetime}
% Added by lyx2lyx
\setlength{\parskip}{\smallskipamount}
\setlength{\parindent}{0pt}

\makeatother

\usepackage{babel}
\providecommand{\exercisename}{Exercise}
\providecommand{\theoremname}{Theorem}

\begin{document}
\renewcommand{\author}{Dr.\,James Ashe}

\newcommand{\classNum}{335}

\newcommand{\classSec}{A}

\newcommand{\nameType}{solutions}

\newcommand{\examType}{Homework 4}

\newcommand{\Date}{\formatdate{6}{2}{2014}} %%\AdvanceDate[12] \today

%\newdate{my_date}{\day}{\month}{\year}%   
%\AdvanceDate[-5] 
%\displaydate{my_date}

%%First page's header%%%%%%%%%%%%%%%%%%%%%%
\fancypagestyle{firststyle} %%header settings for first pagr
{
	\fancyhead[L]{MTH \classNum{}(\classSec{}) \author{}\\
	\textbf{\examType{}}} %\textbf{\examType{}} \Date{}}
	\fancyhead[C]{\textbf{\nameType}}
	\setlength{\headsep}{14pt}
}
%%%%%%%%%%%%%%%%%%%%%%%%%%%%%%%%%%%%%%%%%%

%%Next page's header%%%%%%%%%%%%%%%%%%%%%%%
%\setlist{nolistsep} %eliminates space between list items
\lhead{MTH \classNum{}(\classSec{})} 
\chead{\textbf{\nameType}} 
\cfoot{\thepage{}~of~\pageref{LastPage}} 
\setlength{\headsep}{5pt} 
%%%%%%%%%%%%%%%%%%%%%%%%%%%%%%%%%%%%%%%%%%%

%%Various settings; see the preamble for other settings%%%%%
%%\setenumerate[0]{leftmargin=12pt,itemindent=0pt,labelwidth=0pt}
\renewcommand{\labelenumi}{\textbf{\arabic{enumi}.}}
\renewcommand{\labelenumii}{\textbf{\alph{enumii}.}}
\thispagestyle{firststyle}
%%%%%%%%%%%%%%%%%%%%%%%%%%%%%%%%%%%%%%%%%%

\null

%\vspace{1cm}

\setcounter{section}{3}

\Large{\textbf{Chapter 3}} \large

\textbf{\textbf{\examType} due: \formatdate{27}{3}{2014}}

All non-starred exercises from sections 3.1 and 3.2 (3.1-3.11). 

\setcounter{thm}{0}

\null

%\vspace{1cm}
\begin{xca}
Prove that If $a|b$ and $b|c$ then $a|c$.

\begin{proof}
$a|b\Rightarrow b=na$ for some $n\in\mathbb{Z}$, and $b|c\Rightarrow mb=m\left(na\right)=c$
for some $m\in\mathbb{Z}$. Thus $a|c$ 
\end{proof}
\end{xca}

\begin{xca}
If $a|b$ and $a|c$ then $a|b\pm c$.

\begin{proof}
Let $a|b$ and $a|c$. So then $b=an$ and $c=am$ for some $n,m\in\mathbb{Z}$.
Thus $b\pm c=an\pm am=a\left(n\pm m\right)$ which implies that $a|b\pm c$.
\end{proof}
\end{xca}

\begin{xca}
$\bigstar$ Prove that there are infinitely many prime numbers. 
\end{xca}

\begin{xca}
Determine the following:
\end{xca}

\begin{enumerate}
\item $\gcd(2^{4}3^{2}5\cdot7^{2},2\cdot3^{3}7\cdot11)$

\begin{enumerate}[resume]
\item []$=2\cdot3^{2}\cdot7$
\end{enumerate}
\item The lowest common multiple: $\mbox{lcm}(2^{4}3^{2}5\cdot7^{2},2\cdot3^{3}7\cdot11)$.

\begin{enumerate}[resume]
\item []$=2^{4}\cdot3^{3}\cdot5\cdot7^{2}\cdot11$
\end{enumerate}
\end{enumerate}

\begin{xca}
$\bigstar$ Suppose that $a\neq0$, $b\neq0$, and $\mbox{lcm}(a,b)=m$.
Prove that if $n$ is a common multiple of $a$ and $b$ then $m|n$. 
\end{xca}

\begin{xca}
$\bigstar$ Let $a\neq0$, $b\neq0$, and $d=\gcd\left(a,b\right)$.
Prove that if $d'$ is any common divisor of $a$ and $b$ then $d'|d$.
\end{xca}

\begin{xca}
Find integers $s$ and $t$ so that $1=7\cdot s+11\cdot t$. Show
that $s$ and $t$ are not unique. 

\begin{proof}
Consider $s=-3$ and $t=2$. These can be found as follows:\begin{multicols}{2}

$11=7\cdot1+4$

$7=4\cdot1+3$

$4=3\cdot1+1$

$3=1\cdot3+0$ so $1=\gcd(11,7)$\columnbreak

$1=3\cdot(-1)+4$

$1=\left[7-4\cdot1\right](-1)+4=7(-1)+4(2)$

$1=7(-1)+\left[11-7\right](2)$

$1=7(-3)+11(2)$

\end{multicols}

To see that $s$ and $t$ are unique, consider $s=8$ and $t=-5$.
Note that $as+bt=as+\left[abk-abk\right]+bt=a(s+bk)+b(t-ak)$ for
any $k\in\mathbb{Z}$. Using the answer above, this gives us an infinite
class of solutions. For $k=1$, $7(-3+11)+11(2-7)=7(8)+11(-5)=1$.
\end{proof}
\end{xca}

Equations of the form $ax+by=d$ are called Diophantine equations
(after Diophantus of Alexandria 200-284 AD). Diophantine equations
have a solution if and only if the $\gcd(a,b)|d$, and if a solution
exists then there are infinitely many of them.\vfill

\begin{xca}
Let $d=\gcd(a,b)$. If $a=da'$ and $b=db'$, show that $\gcd(a',b')=1$.

\begin{proof}
By the GCD is a linear combination theorem, $d=ax+by$ for some $x,y\in\mathbb{Z}$.
This implies that $d=da'x+db'y$, and so dividing by $d$ gives $1=a'x+b'y$.
Thus $1=\gcd(a',b')$
\end{proof}
\end{xca}

\begin{xca}
$\bigstar$ If $c^{2}=ab$ and $\gcd\left(a,b\right)=1$, prove that
$a$ and $b$ are perfect squares (Hint: By the theorem above, $a$
and $b$ can be written as the product of primes. They are perfect
squares if and only if each prime in this prime factorization has
an even power). 
\end{xca}

\begin{xca}
Determine $51\mod 13$, $342\mod 85$, and $(82\cdot73)\mod 7$.

\begin{itemize}
\item []$51=3\cdot13+12$ so $12\equiv51\mod 13$
\item []$342=4\cdot85+2$ so $2\equiv342\mod 85$
\item []$82\cdot73=5986=855\cdot7+1$ so $1\equiv(82\cdot73)\mod 7$
\end{itemize}
\end{xca}

\begin{xca}
Let $n$ be a positive integer greater than $1$ and let $a$ and
$b$ be positive integers. Prove that $a\mod n\equiv b\mod n$ if
and only if $a\equiv b\mod n$. 

To prove a ``\textbf{A} if and only if \textbf{B}'' statement, you
must show that \textbf{A} implies \textbf{B }and \textbf{B }implies
\textbf{A }(``if and only if''\textbf{ }is also written ``iff.''
and $\iff$).

We discussed two equivalent ways to define $a=b\bmod n$: \\
$a=b\bmod n\iff n|b-a$, and \\
$a=b\bmod n\iff a=qn+b$ for some $q\in\mathbb{Z}$. 
\begin{proof}
$(a\bmod n=b\bmod n\Rightarrow a=b\bmod n)$ If $a\bmod n=b\bmod n$
then $n|\left(a-b\bmod n\right)$ and $n|\left(b-a\bmod n\right)$.
So then $n|\left(\left[a-b\bmod n\right]-\left[b-a\bmod n\right]\right)$,
and since $-b\bmod n+a\bmod n=0$ the two middle terms cancel out
implying that $n|\left(a-b\right)$. Thus $a=b\bmod n$. 

$\left(a=b\bmod n\Rightarrow a\bmod n=b\bmod n\right)$ Suppose that
$a=b\bmod n$ so then $n|b-a$. Since $n|a-a$, $a=a\bmod n$. It
follows that $n|b-a\Rightarrow n|b-a\bmod n\Rightarrow a\bmod n=b\bmod n$.
\end{proof}
(Alternatively) 
\begin{proof}
$(a\bmod n=b\bmod n\Rightarrow a=b\bmod n)$ If $a\bmod n=b\bmod n$
then $a=nq_{1}+r$ and $b=nq_{2}+r$$\Rightarrow a-b=n(q_{1}-q_{2})\Rightarrow n|a-b\Rightarrow a=b\bmod n$. 

$\left(a=b\bmod n\Rightarrow a\bmod n=b\bmod n\right)$ Suppose that
$a=b\bmod n$ then $a=q_{1}n+r_{1}$ and $b=q_{2}n+r_{2}$. Since
$a=b\bmod n\Rightarrow$$n|a-b\Rightarrow$$n|n\left(q_{1}-q_{2}\right)+\left(r_{1}+r_{2}\right)$.
Thus $n$ divides $n\left(q_{1}-q_{2}\right)$and also divides $r_{1}+r_{2}$.
Since $0\leq r_{1},r_{2}<n$ it follows that $r_{1}-r_{2}=0\Rightarrow r_{1}=r_{2}$.
Thus $a\bmod n=b\bmod n$.\vfill
\end{proof}
\end{xca}


\end{document}
